\documentclass[11pt]{article}
\usepackage{fontspec}
\usepackage[french,english]{babel}
\usepackage{amsmath}
\usepackage{amssymb}
\usepackage{amsthm}
\usepackage{mathrsfs}
\usepackage{tikz-cd}
\usepackage{enumitem}
\usepackage{lmodern}

\newcommand{\printedpage}[1]{\marginpar{\footnotesize p.~#1}}
\newcommand{\msp}[1]{\marginpar{\footnotesize [#1]}}
\newcommand{\hot}{\mathrm{hot}}
\newcommand{\HOT}{\mathrm{HOT}}
\newcommand{\Cat}{\mathrm{Cat}}
\newcommand{\Liss}{\underline{\mathrm{Liss}}}
\newcommand{\multhot}{\mathrm{multhot}}
\newcommand{\mult}{\mathrm{mult}}

\begin{document}

\printedpage{94}
% head: A. Grothendieck : Dérivateurs, Ch. VI, Édition des Universités de Montpellier II et Paris VII
If e.g. $s$ is the final element $e$ of $S$, we find
\[\Phi(\underset{S}{e}) \simeq S,\]
hence
\[\hot_{S,W}(X)\times_S \hot_{S,W}(e)=\hot_{S,W}(X),\]
on the other hand
\[\hot_{S,W}(X\times_S e)=\hot_{S,W}(X_e).\]

\msp{154}
If the two were isomorphic in $\HOT_W(S)$ by the canonical arrow
\[X_e\longrightarrow X,\]
applying the functor $p_!$ (where $p : S\longrightarrow e$), it would follow that $X_e\longrightarrow X$ is in $W$. If for example $S=\mathbf{\Delta}^1$, this would imply that for every $X$ smooth over $S$, $X_1\longrightarrow X$ is in $W$, whereas normally smoothness implies that $X_0\longrightarrow X$ is in $W$. So that would amount to saying that $\Liss\,\mathbf{\Delta}^1=\Liss_0\mathbf{\Delta}^1$, whence at once $\Liss\,S=\Liss_0S$, for every $S$ in Cat, we'd know about it! The most trivial counterexample is to take $X=\{0\}\hookrightarrow S=\mathbf{\Delta}^1$, the inclusion of the ``generic'' open, an open immersion, hence smooth, the fiber $X_1$ is empty! % [sic] s'ensuiverait -> s'ensuivrait, nonstandard verb form

% marginal mark: handwritten sign in the left margin beside the following statement
\textbf{Corollary 1 corrected.} \textit{If $X$ is in} $\Cat/S$\textit{, not necessarily smooth over $S$, then the functor}
\[\mult_{X/S} : Y\longmapsto Y\times_S X : \Cat/S\longrightarrow \Cat/S\]
\textit{induces a functor}
\[\Liss\,S\longrightarrow \Cat/S\]
\textit{compatible with the localizers $W_S^{\mathrm{d}}$, $W_S^{\mathrm{d}}$, hence induces a functor}
\[(\Liss\,S)(W_S^{\mathrm{d}})^{-1}\longrightarrow(\Cat/S)(W_S^{\mathrm{d}})^{-1},\]

\msp{155}
\textit{or again (taking into account the equivalence induced by the inclusion}
\[(\underline{\mathrm{Liss}}\,S)(W_S^{\mathrm{d}})^{-1}\xrightarrow{\ \sim\ }(\Cat/S)(W_S^{\mathrm{d}})^{-1}\quad\textit{),}\] % doubt resolved in DER_VI-c_corrections.csv; input note: underline of Liss -- double underline or single underline
\textit{a functor, again denoted} $\multhot_{X/S}$,
\[\multhot_{X/S} : \HOT_W(S)\longrightarrow\HOT_W(S).\]

This functor therefore satisfies the relation
\[\multhot_{X/S}(\hot_{S,W}(Y))\qquad\underbrace{\ \simeq\ }_{\text{canonical isomorphism}}\qquad\hot_{S,W}(X\times_S Y)\]
% foot: Version 9 novembre 2001
\printedpage{95}
% head: A. Grothendieck : Dérivateurs, Ch. VI, Édition des Universités de Montpellier II et Paris VII
% marginal mark: handwritten sign in the left margin beside the following two lines
for $Y$ \textit{smooth over $S$,} and for arbitrary $Y$ in $\Cat/S$,
\[\multhot_{X/S}(\hot_{S,W}(Y))\qquad\underbrace{\ \simeq\ }_{\text{canonical isomorphism}}\qquad\hot_{S,W}(X\times_S \Phi_S(Y)).\]

When $X$ is itself smooth over $S$, this functor is none other than multiplication by $\hot_{S,W}(X)$:
\[\multhot_{X/S}=\mult_{\hot_{S,W}(X)}\qquad\text{if $X$ \textit{smooth over $S$.}}\]

% marginal mark: handwritten sign in the left margin beside the following paragraph
But its value for $\hot_{S,W}(Y)$, when $Y$ is \textit{not} smooth over $S$, must not be confused with $\hot_{S,W}(X\times_S Y)$.

Thus, if $X$ and $Y$ are in $\Cat/S$, let us consider
\[X\times_S Y\qquad\text{and}\qquad\boxed{\hot_{S,W}(X\times_S Y).}\] % the frame is printed with rounded corners

\begin{enumerate}
\item[a)] If neither $X$ nor $Y$ is assumed smooth over $S$, this element ``has no natural interpretation in terms of operations on the categories $\HOT_W(S)$'' (\footnote[74]{Except in the case where one of the two is in $\underline{\mathrm{Fibhot}}\,S$, cf. annotation on the next page.}).

\msp{156}

\item[b)] If $X$ and $Y$ are smooth over $S$, this element is interpreted as
\[\hot_{S,W}(X)\times_S\hot_{S,W}(Y)\]
(\footnote[75]{But also if \textit{one} of the two elements is in $\underline{\mathrm{Fibhot}}\,S$ (and assuming the axiom L6 bis on $W$).}).

\item[c)] If one of the two $X$, $Y$ is smooth over $S$, without assuming anything about the other, we have an interpretation in terms of $\mult_X$ (if $Y$ smooth), or $\mult_Y$ (if $X$ smooth), in particular
\[\begin{aligned}
\hot_{S,W}(X\times_S Y)\ &\simeq\ \multhot_X(\hot_{S,W}(Y))\quad\text{if $Y$ smooth over $S$}\\
&\simeq\ \multhot_Y(\hot_{S,W}(X))\quad\text{if $X$ smooth over $S$.}
\end{aligned}\]
\end{enumerate}

\textbf{Corollary 2.} \textit{Let $f : S'\longrightarrow S$ be an arrow in} Cat\textit{, whence $f_!$, $f^*$ between} $\HOT_W(S')$ \textit{and} $\HOT_W(S)$\textit{. We have the relation}
\[f_!f^*(\xi)\simeq\multhot_{S'}(\xi),\]
\textit{hence if $S'$ is smooth over $S$,}
\[f_!f^*(\xi)\simeq\hot_{S,W}(S')\times_S\xi.\]

We are now in a position to verify the axioms Der 4$'$b, Der 5$'$ relating to direct images (I, p.~67, 68), and moreover Der 3b (I, p.~111). The axiom Der 4$'$ is none other than the existence of the left adjoints $f_!$ of the $f^*$. The axiom Der 5$'$ allows the ``calculation'' of the fibers of an object $f_!(\xi)$. It says that a certain
% foot: Version 9 novembre 2001
\printedpage{96}
% head: A. Grothendieck : Dérivateurs, Ch. VI, Édition des Universités de Montpellier II et Paris VII
\msp{157}
base change arrow, relative to the commutative diagram in Cat
\[\begin{tikzcd}[row sep=large]
S' \arrow[d,"f"'] & s\backslash S' \arrow[l,"q"'] \arrow[d,"f'"] \\
S & e\simeq\{s\} \arrow[l,"p"']
\end{tikzcd}\qquad(s\in\mathrm{Ob}\,S),\]
namely
\[p^*(f_!(\xi'))\longrightarrow f'_!(q^*(\xi'))\]
is an isomorphism (for $\xi'\in\mathbf{D}(S')$). So let
\[\xi'\in\mathrm{Ob}\,\HOT_W(S'),\qquad\xi'=\hot_{S',W}(X')\]
with $X'$ in $\Cat/S'$. Since
\[f_!(\xi')=\hot_{S,W}(X'),\]
and since $X'$ has no reason to be smooth over $S$ (if we do not assume $S'$ itself smooth over $S$, in which case it might be useful to take $X'$ smooth over $S'$, hence also over $S$), to calculate $p^*(f_!(\xi'))$ there is no question of taking the fiber $X'_s$, but we must take $s\backslash X'=\Phi_S(X')_s$:
\[p^*(f_!(\xi'))\simeq\hot_W(s\backslash X').\]
\[\begin{tikzcd}[row sep=large]
X' \arrow[d] & s\backslash X' \arrow[l] \arrow[d] \\
S' \arrow[d] & s\backslash S' \arrow[l,"q"'] \arrow[d] \\
S & s\backslash S \arrow[l]
\end{tikzcd}\]

Consider the base change diagram along $s\backslash S\longrightarrow S$ opposite \texttt{[i.e. above]}, with cartesian squares. We first see that $q^*(\xi')$ is none other than the class in $\HOT_W(s\backslash S')$ of the inverse image of $X'$ by the base change $s\backslash S'\longrightarrow S'$. To ensure this, it would suffice to take $X'$ smooth over $S'$ (which
\msp{158}
is permissible) - then the description of the inverse image is OK for every base change. But we have seen that in the case of a local co-isomorphism such as $s\backslash S'\longrightarrow S'$, there is not even any need for this artifice - the formula is OK without any smoothness restriction. Thus
\[q^*(\xi')=\hot_{s\backslash S',W}(s\backslash X'),\]
% foot: Version 9 novembre 2001
\printedpage{97}
% head: A. Grothendieck : Dérivateurs, Ch. VI, Édition des Universités de Montpellier II et Paris VII
whence
\[f'_!(q^*(\xi'))\simeq\hot_{S,W}(s\backslash X')=\hot_W(s\backslash X').\]
So we have indeed found
\[p^*f_!\simeq f'_!q^*.\]

We can now also prove Der 3b. We have seen that once we have the left adjoints $f_!$ of the $f^*$, this axiom is stated in \textit{homological form} as follows: If
\[f : S'\longrightarrow S\]
is in $W$, we must prove that for every
\[\xi\in\HOT_W(e)\overset{\text{def}}{=}\mathrm{Hot}_W,\]
the corresponding arrow
\[(*)\qquad p'_!p'^*(\xi)\longrightarrow p_!p^*(\xi)\]
is \texttt{[an]} isomorphism in $\mathrm{Hot}_W$, where $p : S\longrightarrow e$, $p' : S'\longrightarrow e$ are the structural arrows. Now since every object of Cat is smooth over $e$, the preceding corollary 2 allows us to identify this arrow $(*)$ with
\[\eta'\times_S\xi\longrightarrow\eta\times_S\xi\]
\msp{159}
deduced from
\[\hot_W(f) :\quad\underbrace{\ \eta'\ }_{\overset{\text{def}}{=}\,\hot_W(S')}\quad\longrightarrow\quad\underbrace{\ \eta\ }_{\overset{\text{def}}{=}\,\hot_W(S)}\quad.\]
Since the latter is \texttt{[an]} isomorphism by hypothesis, we win.

So we have proved

\textbf{Theorem 3.} \textit{The prederivator} HOT \textit{satisfies the axioms} Der 1, Der 2 (\footnote[76]{\textit{not proved}, the ``result'' I had in mind is botched, cf. supplement proposition 5, page 164.}) \textit{(for the record), and} Der 3b (\footnote[77]{And even the strongest variant Der 3b$'_2$ of Der 3, cf. proof on the next page. \textbf{NB} For Der 3b$'_2$, we used the axiom L6 on $W$ - and this is essential. We see that the validity of Der 3b$'_2$ implies L6, and even the stronger form L6 bis, and even if we assume only Der 3b$'_1$, i.e. $f_0^*$ fully faithful.}), Der 4$'$, Der 5$'$\textit{.} % [sic] canulé (footnote 76) -> annulé (?), not a French word

We also feel like proving the strengthened variants of Der 3, the strongest form being this one (Der 3b$'_2$): if $f\in W$, then $f^*$ induces an \textit{equivalence} of categories
\[f_0^* : \HOT^{\mathrm{lc}}_W(S)\longrightarrow\HOT^{\mathrm{lc}}_W(S').\]

The full faithfulness of the functor $f_0^*$ must follow from the natural variant Der 3b$''_{\mathrm{coh}}$ of Der 3b, once we have developed the formalism of internal Hom in the $\HOT_W(S)$. But to give a meaning to Der 3b$''_{\mathrm{coh}}$, we need the functors \textit{cohomology,} i.e. $p_*$ and $p'_*$ (instead of $p_!$, $p'_!$) for HOT, so this is still premature for the moment. The construction of the $f_*$ and that of internal Hom will moreover surely be done $\pm$ jointly.
% foot: Version 9 novembre 2001

\printedpage{98}
% head: A. Grothendieck : Dérivateurs, Ch. VI, Édition des Universités de Montpellier II et Paris VII
\msp{160}
But I realize just now that we find almost ``by eye'' the fact that $f_0^*$ is indeed an \textit{equivalence} of categories, and that $f_!^0$ is the quasi-inverse equivalence, by using the adjunction of the functors $f_!^0$, $f_0^*$, and recalling that to prove that they are equivalences quasi-inverse to one another, it amounts to the same thing to show that the adjunction morphisms
\[\begin{aligned}
\mathrm{id}&\quad\xrightarrow{\ \alpha_0\ }\quad f_0^*f_!^0\\
\mathrm{id}&\quad\xleftarrow{\ \beta_0\ }\quad f_!^0f_0^*
\end{aligned}\]
are \texttt{[some]} isomorphisms. I am going to try to describe these adjunction arrows in the general case (without assuming $f\in W$), then to show that they are in $W$ if $f$ is.

\begin{enumerate}
\item[\texttt{[1)}] \texttt{Let]} $\xi=\hot^{\mathrm{lc}}_{S,W}(X)$ in $\HOT_W(S)$. We may assume $X\in\mathrm{Ob}\,\underline{\mathrm{Parf}}\,S$. Since \texttt{[``Since'' crossed out]} $X$ smooth over $S$, we have
\[f_0^*(\xi)=\hot^{\mathrm{lc}}_{S',W}(X\times_S S'),\qquad f_!^0f_0^*(\xi)=\hot^{\mathrm{lc}}_{S,W}(X\times_S S').\]
So
\[\beta_{0,\xi} : f_!^0f_0^*(\xi)\longrightarrow\xi\]
is an arrow
\[X\times_S S'\longrightarrow X\qquad\text{in }(\Cat/S)(W_S)^{-1}\]
and I presume that it is $\mathrm{pr}_1$ (\footnote[78]{\textbf{NB} We see that the axiom L6 is $\pm$ \textit{necessary} for the result on $f_0^*$, $f_!^0$ to be valid.}). Is it in $W$ if $f : S'\longrightarrow S$ is? Yes (for $X$ is perfect

\msp{161}

over $S$), provided that $W$ satisfies the axiom L6 (cf. page 43).

\item[2)] Let
\[\xi'=\hot^{\mathrm{lc}}_{S',W}(X')\qquad\text{in }\HOT_W(S'),\]
we have
\[f_!^0(\xi')=\hot^{\mathrm{lc}}_{S,W}(X'),\]
but to make $f_0^*$ of this element explicit, since $X'$ has no reason to be Hot-fibered over $S$ (in which case it would suffice to take $f^*(X')\overset{\text{def}}{=}X\times_S S'$), we must first take $C_S(X')$, then apply $f^*$:
\[f_0^*f_!^0(X')=\hot^{\mathrm{lc}}_{S',W}(C_S(X')\times_S S').\]

We must therefore define
\[X'\xrightarrow{\beta_{0,X'}}C_S(X')\times_S S'\qquad\text{in }(\Cat/S')W_{S'}^{-1}.\]

Now we find a canonical arrow in $\Cat/S'$, deduced from the $S$-morphism $i_{X'}$: $X'\longrightarrow C_S(X')$, i.e. from the commutative square in solid lines
\end{enumerate}
% foot: Version 9 novembre 2001

\end{document}
